\chapter{\rm\bfseries Results}
\label{ch:results}

This chapter reports the empirical findings across all six gripper
variants, namely \path{finray_18}, \path{finray_18_model2},
\path{finray_26}, \path{finray_26_5deg},
\path{parallel_jaw_stock}, and \path{parallel_jaw_TPU}. The findings
are organised around the primary measurement, the contact force $F_z$,
at swept command depths and $\pm 2.5^{\circ}$ tilts about the $x$ and
$y$ axes. All values
were $F_z$-bias-corrected using the per-session Franka $F_{\mathrm{ext}}$ zero-load
reading. The corrected values were then pooled across trials with
matching command depths and tilts unless otherwise noted.
Interpretation is deferred to Chapter~\ref{ch:discussion}.

\section{Print-to-Print Manufacturing Variability}
\label{sec:manufacturing-results}

Figure~\ref{fig:m1-vs-m2} compares two independent prints of the same
\texttt{finray\_18} CAD geometry on the calibrated depth axis. Each
print was zeroed at its own first majority-success cell
(Section~\ref{sec:dcal}). The two prints differed by a $+4.6$\,N mean
force offset, with the second print being the stiffer, and the RMSE
over the 10\,mm window was $5.4$\,N. Removing the constant offset left
a residual RMSE of $2.9$\,N. This $\approx 5$\,N offset was adopted as
the canonical same-session print-to-print variability, approximately
$15\%$ of the knee peak force.

Cross-session drift on the \emph{same} print was measured by comparing
\texttt{finray\_18} data collected on separate days approximately three
weeks apart. That drift was approximately $7$--$10$\,N at knee and
knee-adjacent depths. Therefore, drift across sessions on a single
specimen was larger than same-session print-to-print variability.

\begin{figure}[ht]
\centering
\includegraphics[width=0.85\linewidth]{fr18_m1_vs_m2_aligned0.png}
\caption[Print-to-print force comparison]{Contact force versus
calibrated depth for two independently printed copies of
\texttt{finray\_18}, zero tilt, successful cells only. Error bars: one
standard deviation across repeats; where no bar is visible, the spread
is smaller than the marker.}
\label{fig:m1-vs-m2}
\end{figure}

The two prints' primary force-depth response was otherwise nearly
identical (pairwise shape NMSE $0.41\%$, Table~\ref{tab:shape-nmse}).
Their secondary x-axis bias also shared the same depth-gated envelope
as the other bias measurements in this work
(Section~\ref{sec:bipolar-bias-results}), being largest near the
boundary and relaxing toward zero at depth. The two prints differed
only in where that shared envelope sat relative to zero. The bias of
\texttt{finray\_18} crossed zero, running negative at the boundary,
briefly positive through the knee, and negative again at depth.
However, the bias of \texttt{finray\_18\_model2} stayed negative
throughout the full measured range. The gap between the two was itself
depth-gated, at $2$--$3.5$\,N near the boundary and shrinking to
$<1$\,N by $d=+10$\,mm. The attribution of this offset is
deferred to Section~\ref{sec:error-discussion}.

\section{Working Range at Zero Tilt}
\label{sec:working-range-zero}

Figure~\ref{fig:working-range-hist} displays the successful-grasp
working range of each variant at zero tilt, on the calibrated depth
axis of Section~\ref{sec:dcal}. The parallel-jaw variants exhibited
closed windows of $1.5$\,mm (\texttt{parallel\_jaw\_stock}) and
$2.5$\,mm (\texttt{parallel\_jaw\_TPU}). In contrast, all four finray
variants spanned the full $\geq 10$\,mm analysis window. Their ranges
were right-censored at the window end, and the underlying record
extends beyond it.

\begin{figure}[ht]
\centering
\includegraphics[width=0.85\linewidth]{working_range_hist.png}
\caption[Zero-tilt working range per variant]{Successful-grasp working
range on the calibrated depth axis, zero tilt. Hatched bars are
right-censored at the $10$\,mm analysis-window end (Section~\ref{sec:dcal}).}
\label{fig:working-range-hist}
\end{figure}

\section{Working Range Under Tilt}
\label{sec:pj-vs-finray-results}

Figure~\ref{fig:working-range-tilt} displays the same working ranges
over every tested orientation. Extending the comparison in this way
narrowed the parallel-jaw windows further, to $0.5$--$2.5$\,mm. One
stock-jaw tilt condition was reduced to a single successful depth
($0$\,mm span), and two others produced no successful grasp at any
depth. Five orientations were tested per finray variant and four per
parallel-jaw variant, whose $\theta_x = -2.5^{\circ}$ condition is
covered by mirror symmetry instead of by separate cells. Every finray
variant retained its full $\geq 10$\,mm window under every orientation
tested.

\begin{figure}[ht]
\centering
\includegraphics[width=\linewidth]{working_range_tilt.png}
\caption[Working range under tilt, all variants]{Successful-grasp
working range on the calibrated depth axis, all tested orientations
per variant (five for the finrays, four for the parallel jaws). \textcolor{red}{$\times$}: no successful attempt at any
depth. Hatched bars are right-censored at the analysis-window end.}
\label{fig:working-range-tilt}
\end{figure}

The parallel jaws auto-failed on the Franka excess-force safety monitor
at depths only $1$--$1.5$\,mm beyond their successful-grasp window.
That monitor triggers at $-60$\,N in $F_z$. However, the finrays never
triggered the safety monitor within the tested $-2$ to $+16$\,mm range
on any variant.

\section{Tilt-Aligned Response Comparison}
\label{sec:tilt-similarity-results}

Figure~\ref{fig:tilt-similarity} displays the five tilt-condition
curves of each finray variant together on the calibrated depth axis.
Within each variant, the orientations tracked one another closely. The
mean spread across the five orientations, averaged over the calibrated
window, was $1.3$--$2.4$\,N depending on variant, which is small
relative to the $31$--$53$\,N knee peaks reported in
Table~\ref{tab:zero-tilt-summary}. The curves separated near the
boundary and converged in the plateau, consistent with the depth-gated
tilt robustness reported in Section~\ref{sec:bipolar-bias-results}.

\begin{figure}[ht]
\centering
\includegraphics[width=0.95\linewidth]{finray_tilt_similarity.png}
\caption[Tilt-aligned response per finray variant]{Contact force versus
calibrated depth for each finray variant, all five tested orientations
overlaid (0\textdegree, $x\pm2.5^{\circ}$, $y\pm2.5^{\circ}$). Within
each panel the five curves nearly coincide.}
\label{fig:tilt-similarity}
\end{figure}

\section{One Shared Curve Shape Across the Finray Family}
\label{sec:shape-results}
\label{sec:zero-tilt-results}

Under zero tilt, all four finray variants exhibited a stereotyped
\emph{bell-plus-valley} force-depth response with four regions. These
were (i)~a boundary regime where contact force rose steeply with press
depth, (ii)~a knee peak at calibrated depths of $+2.5$ to $+4.5$\,mm,
(iii)~a valley trough near $+6$ to $+8.5$\,mm, and (iv)~a recovery
regime toward the window end. Table~\ref{tab:zero-tilt-summary} lists
the key features for each variant. The four features lay within
$\approx \pm 1.5$\,mm of the same depth locations across all four
variants, and the absolute force magnitudes scaled approximately with
the frontal contact area of the finger.

\begin{table}[ht]
\centering
\caption[Zero-tilt curve summary]{Key features of the $0^{\circ}$
force-depth curve for each finray variant, on the calibrated depth axis
(Section~\ref{sec:dcal}). Peak $|F_z|$ is the maximum
observed contact force in the knee region; valley $|F_z|$ is the local
minimum between the knee peak and the deep-regime recovery. The
\texttt{finray\_26} knee is a flat plateau spanning $+2.5$ to
$+4.5$\,mm; its listed peak depth is the maximum-force cell.}
\label{tab:zero-tilt-summary}
\begin{tabular}{lcccc}
\toprule
Variant & Knee peak $|F_z|$ & Peak depth & Valley $|F_z|$ & Valley depth \\
\midrule
\texttt{finray\_18}        & $30.6$\,N & $+3.0$\,mm & $27.1$\,N & $+7.5$\,mm \\
\texttt{finray\_18\_model2} & $36.8$\,N & $+2.5$\,mm & $29.5$\,N & $+6.0$\,mm \\
\texttt{finray\_26}        & $44.3$\,N & $+4.5$\,mm & $35.5$\,N & $+8.5$\,mm \\
\texttt{finray\_26\_5deg}  & $53.1$\,N & $+2.5$\,mm & $39.1$\,N & $+6.5$\,mm \\
\bottomrule
\end{tabular}
\end{table}

Normalising each variant's curve by its own peak force and averaging
over its five tested orientations makes this shared shape directly
visible (Figure~\ref{fig:shape-normalized}): all four finray variants
collapse onto a single rise--peak--valley--recovery curve over the full
$[0,10]$\,mm calibrated window. Table~\ref{tab:shape-nmse} quantifies
the pairwise agreement between variants as a normalised mean-squared
error (NMSE) and, in decibels, $10\log_{10}(\text{NMSE})$; the mean
pairwise NMSE was $0.58\%$ ($-22.4$\,dB), with the largest
disagreement between any two variants at $0.79\%$ ($-21.0$\,dB).

\begin{figure}[ht]
\centering
\includegraphics[width=0.75\linewidth]{finray_shape_normalized.png}
\caption[Shared force-depth shape across the finray family]{Peak-normalised
$|F_z|$ versus calibrated depth $d_{\text{cal}}$, averaged over all five
tested orientations per variant (shaded band: across-orientation spread).
All four finray variants collapse onto one shared response shape.}
\label{fig:shape-normalized}
\end{figure}

\begin{table}[ht]
\centering
\caption[Pairwise shape agreement]{Pairwise shape agreement between
finray variants over the full $[0,10]$\,mm calibrated window: NMSE (top),
NMSE in decibels (middle), and RMSE as a percentage of peak (bottom,
in brackets).}
\label{tab:shape-nmse}
\small
\begin{tabular}{lcccc}
\toprule
 & \texttt{finray\_18} & \texttt{finray\_18\_model2} & \texttt{finray\_26} & \texttt{finray\_26\_5deg} \\
\midrule
\texttt{finray\_18} & --- &
  \begin{tabular}{c}$0.41\%$\\($-23.9$\,dB)\\{}[$5.5\%$]\end{tabular} &
  \begin{tabular}{c}$0.69\%$\\($-21.6$\,dB)\\{}[$7.2\%$]\end{tabular} &
  \begin{tabular}{c}$0.79\%$\\($-21.0$\,dB)\\{}[$7.8\%$]\end{tabular} \\[2ex]
\texttt{finray\_18\_model2} & & --- &
  \begin{tabular}{c}$0.13\%$\\($-28.8$\,dB)\\{}[$3.1\%$]\end{tabular} &
  \begin{tabular}{c}$0.70\%$\\($-21.6$\,dB)\\{}[$7.1\%$]\end{tabular} \\[2ex]
\texttt{finray\_26} & & & --- &
  \begin{tabular}{c}$0.75\%$\\($-21.2$\,dB)\\{}[$7.4\%$]\end{tabular} \\[2ex]
\texttt{finray\_26\_5deg} & & & & --- \\
\bottomrule
\end{tabular}
\end{table}

\section{Bipolar Tilt Bias on the Finray Family}
\label{sec:bipolar-bias-results}

A $\pm 2.5^{\circ}$ tilt produced a signed bias
$b(d,a) = |F_z^{+2.5^{\circ}}(d,a)| - |F_z^{-2.5^{\circ}}(d,a)|$.
Depths in this section and the next are quoted on the raw command-depth
axis $d$, the native axis of the per-cell bias record, rather than the
calibrated axis $d_{\text{cal}}$ used elsewhere in this chapter. The
bias's depth trajectory itself followed a bipolar shape: a strong signed
peak at boundary depths, a zero-crossing near the knee, a sign-reversed
plateau through the valley, and depth-gated attenuation toward zero at
deep press. The clearest example was the y-axis bias on
\texttt{finray\_26\_5deg}, which had full coverage across 14 sign-flip
pairs from $-2$ to $+16$\,mm:

\begin{itemize}
  \item Boundary peak: $b = +5.40 \pm 0.12$\,N at $d = -1.5$\,mm.
  \item Zero crossing: between $d = -0.5$\,mm ($b = +2.00$\,N) and
        $d = +1.5$\,mm ($b = -1.35$\,N).
  \item Knee plateau: $b = -1.57 \pm 0.21$\,N at $d = +2.0$\,mm,
        persisting to $d = +3.0$\,mm.
  \item Depth-gated recovery: $|b| < 0.9$\,N at all depths beyond
        $d = +4$\,mm.
\end{itemize}

The x-axis bias on \texttt{finray\_26\_5deg} had a similar bipolar
shape but with opposite sign at the boundary
($b_x = -4.17 \pm 0.32$\,N at $d = -1.5$\,mm).

\section{Cross-Axis Sign Reversal at the Boundary}
\label{sec:cross-axis-results}

At boundary depths, the x-axis and y-axis biases had \emph{opposite
signs} on \texttt{finray\_26\_5deg} (Table~\ref{tab:cross-axis-sign}).

\begin{table}[ht]
\centering
\caption[Cross-axis bias sign reversal at boundary]{Bias on
\texttt{finray\_26\_5deg} at boundary depths. x-axis and y-axis biases
had opposite signs (sign product $-1$). Values are
$|F_z^{+2.5^{\circ}}|-|F_z^{-2.5^{\circ}}|$ in newtons.}
\label{tab:cross-axis-sign}
\begin{tabular}{lccc}
\toprule
Depth & x-bias & y-bias & Sign product \\
\midrule
$-1.5$\,mm & $-4.17 \pm 0.32$ & $+5.40 \pm 0.12$ & $-1$ \\
$-1.0$\,mm & $-3.94 \pm 0.18$ & $+4.27 \pm 0.07$ & $-1$ \\
$-0.5$\,mm & $-3.51 \pm 0.39$ & $+2.00 \pm 0.24$ & $-1$ \\
\bottomrule
\end{tabular}
\end{table}

The same sign-reversal pattern was present on \texttt{finray\_18} at
$d = -1.0$\,mm ($x$-bias $= -2.44 \pm 0.17$\,N; $y$-bias
$= +4.49 \pm 0.47$\,N), demonstrating that the effect was not specific to
the $5^{\circ}$ mount condition.

In absolute terms, the boundary y-axis asymmetry corresponds to a
$2.36\times$ force ratio between the two tilt signs. On
\texttt{finray\_26} at $d=-1.0$\,mm, the mean stationary force was
$8.00$\,N under $+2.5^{\circ}$ against $3.39$\,N under $-2.5^{\circ}$,
while at $d=+10.0$\,mm the same pair differed by under $1\%$, at
$36.26$\,N against $35.95$\,N. It is worth noting that these values
come from the four cells of the tilt-geometry verification
(Table~\ref{tab:tilt-geometry}), which cover one variant on one axis,
so the ratio should be read as an illustration of the depth gating
reported in Section~\ref{sec:bipolar-bias-results} and not as a
separately established quantity.

\section{Mount-Tilt Alignment: $5^{\circ}$ Mount as a Depth Shift}
\label{sec:mount-tilt-results}

After a two-parameter best-fit alignment, a single depth shift plus a
single force offset, the $5^{\circ}$-mounted variant collapsed onto
the flat-mounted variant with residual RMSE below $2$\,N on both tilt
axes independently (Table~\ref{tab:mount-tilt-alignment}).

\begin{table}[ht]
\centering
\caption[Mount tilt equivalence]{Best-fit alignment of
\texttt{finray\_26\_5deg} onto \texttt{finray\_26}. The two axes yielded
independent estimates of the depth shift within $0.02$\,mm of each other.}
\label{tab:mount-tilt-alignment}
\begin{tabular}{lccc}
\toprule
Axis & Depth shift & Force offset & Alignment RMSE \\
\midrule
x     & $-1.14$\,mm & $+2.7$\,N & $1.75$\,N \\
y     & $-1.12$\,mm & $+3.1$\,N & $1.58$\,N \\
\emph{mean} & $\mathbf{-1.13}$\,\textbf{mm} & $+2.9$\,N & $1.66$\,N \\
\bottomrule
\end{tabular}
\end{table}
